% Einstellungen und Pakete für die IWES Vorlage
% Unterstützt: LuaLaTeX (Rotis-Schrift) und pdfLaTeX (Cabin-Fallback)

%=== Kompatibilität ====================================
\usepackage{scrhack}    % KOMA-Script Kompatibilität

%=== Schriften =========================================
\usepackage{iftex}

\ifluatex
    \usepackage{fontspec}
    \IfFileExists{./font/RotisSansSerifStd-Regular.otf}{%
        \setsansfont{RotisSansSerifStd-Regular.otf}[
            Path           = ./font/,
            BoldFont       = RotisSansSerifStd-Bold.otf,
            ItalicFont     = RotisSansSerifStd-Italic.otf,
            BoldItalicFont = RotisSansSerifStd-Bold.otf,
            SmallCapsFont  = RotisSansSerifStd-Regular.otf,
            SmallCapsFeatures = {Scale=0.8},
        ]
        \renewcommand*\familydefault{\sfdefault}
    }{%
        \setsansfont{TeX Gyre Heros}
        \renewcommand*\familydefault{\sfdefault}
        \PackageWarning{IWES}{Rotis nicht gefunden - verwende TeX Gyre Heros}
    }
\else
    \usepackage[T1]{fontenc}
    \usepackage[utf8]{inputenc}
    \usepackage[sfdefault]{cabin}
\fi

%=== Sprache ===========================================
\usepackage[english, ngerman]{babel}    % Beide geladen; aktive Sprache wird per \selectlanguage in 00_main.tex gesetzt
\usepackage{csquotes}       % Für biblatex erforderlich
\usepackage{etoolbox}       % Für Bedingungs-Makros (\ifdefstring)
\usepackage{hyphenat}

%=== Arbeitstyp-Konfiguration ===========================
% Befehl: \setthesistype{master|bachelor|ip|wai}
% Setzt automatisch alle typ-spezifischen Variablen
\newcommand{\setthesistype}[1]{%
    \def\thesistype{#1}%
    % Master oder Bachelor
    \ifboolexpr{ test {\ifdefstring{\thesistype}{master}} or test {\ifdefstring{\thesistype}{bachelor}} }{%
        \def\thesisdesc{Wissenschaftliche Arbeit zur Erlangung des akademischen Grades}%
        \def\thesismodule{}%
        \newbool{showcoexaminer}\booltrue{showcoexaminer}%
        \newbool{showdegree}\booltrue{showdegree}%
    }{}%
    % Interdisziplinäres Projekt
    \ifdefstring{\thesistype}{ip}{%
        \def\thesisdesc{Ausarbeitung für das Modul}%
        \def\thesismodule{Interdisziplinäres Projekt}%
        \newbool{showcoexaminer}\boolfalse{showcoexaminer}%
        \newbool{showdegree}\boolfalse{showdegree}%
    }{}%
    % Wissenschaftliches Arbeiten im Ingenieurwesen
    \ifdefstring{\thesistype}{wai}{%
        \def\thesisdesc{Ausarbeitung für das Modul}%
        \def\thesismodule{Wissenschaftliches Arbeiten im Ingenieurwesen}%
        \newbool{showcoexaminer}\boolfalse{showcoexaminer}%
        \newbool{showdegree}\boolfalse{showdegree}%
    }{}%
}

%=== Seitenlayout ======================================
\usepackage[
    a4paper,
    includeheadfoot,
    left=2.5cm,
    right=2.5cm,
    top=2cm,
    bottom=0.5cm
]{geometry}

\usepackage{pdfpages}       % PDF-Dateien einbinden

%=== Kopf-/Fußzeilen ===================================
\usepackage[headsepline, automark]{scrlayer-scrpage}
\clearpairofpagestyles
\automark[section]{chapter}

\makeatletter
\if@twoside
    \lehead{\pagemark}
    \rehead{\slshape\leftmark}
    \lohead{\slshape\rightmark}
    \rohead{\pagemark}
\else
    \lohead{\slshape\headmark}
    \rohead{\pagemark}
\fi
\makeatother

\pagestyle{scrheadings}

% Kompatibilitätsmakros: \lhead{}/\rhead{} in 00_main.tex sind Kurzbefehle für
% laufende Kolumnentitel. Es handelt sich NICHT um fancyhdr-Befehle — sie rufen
% intern \markboth/\markright von KOMA-Script auf.
\DeclareDocumentCommand{\lhead}{m}{\markboth{#1}{#1}}
\DeclareDocumentCommand{\rhead}{m}{\markright{#1}}
\DeclareDocumentCommand{\chead}{m}{}

%=== Institutionelle Konstanten ========================
% Diese Werte anpassen, falls sich Institutsname, Professor oder Adresse ändern.
\def\universityName{Leibniz Universität Hannover}
\def\instituteName{Institut für Windenergiesysteme}
\def\professorName{Prof. Dr.-Ing. Andreas Reuter}
\def\instituteStreet{Appelstraße 9A}
\def\instituteCity{30167 Hannover}

%=== Zeilenabstand =====================================
\usepackage{setspace}
\onehalfspacing

%=== Hyperlinks ========================================
\ifluatex
    \usepackage[verbose=silent]{microtype}
\else
    \usepackage{microtype}
\fi

\usepackage{hyperref}
\hypersetup{
    hypertexnames=false,
    colorlinks=true,
    linkcolor=blue!70!black,
    citecolor=green!50!black,
    urlcolor=blue,
    filecolor=magenta!70!black,
}

\usepackage{xurl}           % Besserer URL-Umbruch

%=== Tabellen ==========================================
\usepackage{colortbl}
\usepackage{booktabs}       % Professionelle Tabellen
\usepackage{tabularx}
\usepackage{longtable}
\usepackage{array}
\newcolumntype{P}[1]{>{\centering\arraybackslash}p{#1}}
\usepackage{rotating}
\usepackage{multirow}
\usepackage{multicol}

%=== KI-Dokumentation ==================================
% Spaltentyp und Boxen für die Dokumentation des KI-Einsatzes (91_ki_dokumentation.tex)
\newcolumntype{L}[1]{>{\raggedright\arraybackslash}p{#1}}
\usepackage{tcolorbox}
\tcbuselibrary{breakable, skins}
\definecolor{kiakzent}{rgb}{0.0,0.314,0.608}   % LUH-Blau
% Prompt-Box:   \begin{kiprompt}[Titel] ... \end{kiprompt}
\newtcolorbox{kiprompt}[1][Prompt]{%
    enhanced, breakable, sharp corners,
    colback=kiakzent!6, colframe=kiakzent, boxrule=0pt, leftrule=2.5pt,
    left=6pt, right=6pt, top=4pt, bottom=4pt,
    fonttitle=\bfseries\footnotesize, coltitle=kiakzent, colbacktitle=kiakzent!6,
    title={#1}, attach title to upper={\par\smallskip}, fontupper=\small,
    before skip=6pt, after skip=4pt,
}
% Antwort-Box:  \begin{kiantwort}[Titel] ... \end{kiantwort}
\newtcolorbox{kiantwort}[1][Antwort]{%
    enhanced, breakable, sharp corners,
    colback=black!3, colframe=black!45, boxrule=0pt, leftrule=2.5pt,
    left=6pt, right=6pt, top=4pt, bottom=4pt,
    fonttitle=\bfseries\footnotesize, coltitle=black!60,
    title={#1}, attach title to upper={\par\smallskip}, fontupper=\small,
    before skip=4pt, after skip=10pt,
}

%=== Grafiken ==========================================
\usepackage{graphicx}
\graphicspath{{figures/}}   % Bilderpfad
\usepackage{float}
\usepackage{wrapfig}
\usepackage{caption}
\usepackage{subcaption}
\usepackage[section]{placeins}

%=== Mathematik ========================================
\usepackage{amsmath}
\usepackage{amsfonts}
\usepackage{amssymb}

%=== Listen ============================================
\usepackage{enumitem}       % Flexible Listenabstände (ersetzt paralist)
\setlist{nosep}             % Kompakte Listen global (kein zusätzlicher vertikaler Abstand)

%=== Einheiten und Symbole =============================
\usepackage{siunitx}
\sisetup{locale=DE}         % Deutsche Formatierung
\usepackage[right]{eurosym}

%=== Typographie =======================================
\clubpenalty = 10000
\widowpenalty = 10000
\displaywidowpenalty = 10000
\tolerance = 10000
\hbadness = 10000
\emergencystretch = 3em
\hfuzz = 1pt

%=== Inhaltsverzeichnis ================================
\RedeclareSectionCommand[tocpagenumberwidth=2.5em]{chapter}

%=== Literaturverzeichnis (biblatex) ===================
\usepackage[
    backend=biber,
    style=numeric,
    sorting=none,
    doi=true,
    url=true,
    isbn=true,
    maxbibnames=99,
    giveninits=true,
]{biblatex}

\addbibresource{99_bibliography.bib}

% DOI als klickbaren Link formatieren
\DeclareFieldFormat{doi}{%
    \mkbibacro{DOI}\addcolon\space
    \ifhyperref
        {\href{https://doi.org/#1}{\nolinkurl{#1}}}
        {\nolinkurl{#1}}%
}

% Warnungen bei fehlenden DOI/URL/ISBN
\AtEveryBibitem{%
    % Artikel: DOI oder URL erforderlich
    \ifentrytype{article}{%
        \iffieldundef{doi}{%
            \iffieldundef{url}{%
                \PackageWarning{biblatex}{Fehlende DOI/URL für Artikel: \thefield{entrykey}}%
            }{}%
        }{}%
    }{}%
    % Bücher: ISBN oder DOI erforderlich
    \ifentrytype{book}{%
        \iffieldundef{isbn}{%
            \iffieldundef{doi}{%
                \PackageWarning{biblatex}{Fehlende ISBN/DOI für Buch: \thefield{entrykey}}%
            }{}%
        }{}%
    }{}%
}
